\documentclass[11pt,a4paper]{article}

\usepackage[top=25mm,bottom=25mm,left=25mm,right=25mm]{geometry}
\usepackage[T1]{fontenc}
\usepackage[utf8]{inputenc}
\usepackage{mathpazo}
\usepackage{enumitem}
\usepackage{titlesec}
\usepackage{parskip}

\setlength{\parskip}{0.4em}
\setlength{\parindent}{0pt}

\titleformat{\section}
  {\normalfont\bfseries\large}
  {\thesection.}{0.5em}{}
\titlespacing*{\section}{0pt}{1.8em}{0.6em}

\pagestyle{empty}

\begin{document}

{\Large\bfseries Q\&A with Irina Scherbakova and\\ Viktoriya Sukovata on \textit{The Commissar}}

\smallskip

{\small Kino Kaddish screening, Karlshorst, 2026}

\bigskip

\section*{I.\quad Grossman, Berdychiv, and the Choice of Source}

The film goes back to Vasily Grossman's short story ``In the Town of Berdichev,'' and it keeps Berdychiv as its setting. Berdychiv was not just any provincial town. It was one of the most important Jewish cities in the Russian Empire: dense Jewish life, multilingualism, and later, extreme violence. What drew Askoldov to this story, and what does Berdychiv bring to the film as a setting?

\section*{II.\quad Pogrom Violence and Holocaust Foreshadowing}

One of the most striking moments in the film is the visionary sequence sometimes called the ``Passage of the Doomed.'' We see Civil War pogrom imagery shade into something unmistakably connected to the Holocaust: striped uniforms, yellow stars, smoke. This sequence does not exist in Grossman's story. It is entirely Askoldov's invention. How do you read that scene? Does it ask us to see the Civil War through the Holocaust, or is Askoldov pointing to a longer continuity of anti-Jewish violence? And more broadly, can we understand the film as a meditation on the fate of Jews in the twentieth century, told through the lens of pogroms, imperial collapse, and civil war?

\section*{III.\quad The Magazanik Family and Jewish Representation}

At the center of this film is not a heroic military collective but the Magazanik family: their cramped domestic space, their humor, their fear, their way of surviving. In Soviet cinema, Jewish characters were usually marginal or caricatured. Here they are at the moral and emotional core. At the same time, some scholars have argued that the film still draws on inherited Russian literary and Christianized stereotypes. How do you assess the film's representation of Jewish life? What is genuinely new about it, and where does it remain constrained by older conventions?

\section*{IV.\quad Vavilova, Gender, and Revolutionary Identity}

The film is called \textit{The Commissar}, and Vavilova's pregnancy and transformation are at the center of its structure. She arrives as a figure of revolutionary authority, and then her encounter with the Magazanik family forces a confrontation between ideological abstraction and embodied, domestic, Jewish life. How does that encounter reshape what the film says about revolutionary identity? And what role do gender, maternity, and the body play in that shift?

\section*{V.\quad Genre, Canon, and Censorship}

The film is set during the Civil War, not the Great Patriotic War, but it still enters into dialogue with the whole tradition of Soviet war cinema: heroism, sacrifice, ideological clarity. To what extent does \textit{The Commissar} belong to that tradition, and to what extent does it work against it?

\medskip

The film was finished in 1967 and then banned for twenty years. That suggests it touched several nerves at once: Jewish visibility, pogrom violence, the dismantling of Civil War heroism, moral ambiguity. What was most threatening to the Soviet authorities? Was there one element that made the film unacceptable, or was it the combination?

\section*{VI.\quad Afterlife, Reception, and the Film Today}

After 1991, the film entered different memory debates in Ukraine and Russia. How has the reception of \textit{The Commissar} changed in those two contexts? What has been emphasized, and what has been neglected or reframed?

\medskip

The history of this film is also the history of its return: to audiences, to critics, to festivals, and eventually to cultural canons. What were the key moments in that process of rehabilitation, and how secure is the film's place today?

\medskip

Finally: we are screening this film tonight in Berlin, in 2026, for an audience shaped in part by post-Soviet Jewish migration, including families with roots in Ukraine who still often move within Russian-speaking cultural worlds. What does \textit{The Commissar} open up for such an audience today? What can this film make visible, or make discussable, in a way that other frameworks of memory cannot?

\end{document}
